\chapter{Conclusion}
\label{ch:conclusion}

This thesis asked whether AI services could be provisioned at the network level, discovered automatically and consumed without credentials or configuration. Saturn answers that question with a protocol, three implementations, and an analytical evaluation. The three implementations tested the protocol in different settings: a VLC media player extension brought AI into a consumer application that has never had it, a router firmware integration provisioned AI on a router, and an AI coding agent fork handled streaming tool calls and multi-turn conversation. The specification, not a reference implementation, carries the interoperability. Saturn's current reach is bounded by institutional WiFi's AP isolation, which blocks multicast discovery on the networks where access equity matters most. A hybrid architecture pairing mDNS with a unicast HTTPS fallback would extend Saturn to these environments, and controlled user studies would test whether the analytical step reductions translate to real usability gains. The gap between AI's collapsing inference costs and its persistent access barriers is a provisioning problem. Saturn demonstrates that existing protocols mDNS/DNS-SD can apply to AI services in order to make them more accessibile and require less configuration then current alternatives

